\documentclass[11pt]{article}


\setlength{\parindent}{0pt}
\usepackage{xltxtra}
\usepackage{hyperref}
\usepackage{tabularx}
\usepackage{enumitem}
\hypersetup{hidelinks}
\usepackage{url}
\urlstyle{tt}
\usepackage{xcolor}
\usepackage{amsmath}
\definecolor{CVBlue}{RGB}{23,110,191}
\usepackage{calc}
\usepackage{graphicx}
\usepackage{tikz}
\usetikzlibrary{calc}
\usepackage{fontspec}
\usepackage{xeCJK}
\usepackage{enumitem}
\CJKsetecglue{} %% 取消中文与数字之间的间隙


%% ================== 彻底修复后的中英文主体字体配置 ==================
%% 1. 主文档西文字体设置（主导英文、数字、论文标题斜体）
\setmainfont[ 
    Path = fonts/Main/, 
    Extension = .otf,
    BoldFont = texgyretermes-bold.otf,
    ItalicFont = texgyretermes-italic.otf,       % 完美调用本地原装斜体
    BoldItalicFont = texgyretermes-bolditalic.otf % 完美调用本地原装粗斜体
]{texgyretermes-regular.otf} % 正文字体

%% 2. 中文字体设置（主导简历中的所有汉字，绝不乱入英文配置）
\setCJKmainfont[ 
    Path = fonts/hansans/, 
    Extension = .ttf, 
    BoldFont = NotoSansSC-Bold.ttf,
    AutoFakeSlant = 0.2 % 中文若遇斜体，自动由算法转换为伪斜体
]{NotoSansSC-Regular.otf}
%% ==================================================================

\usepackage{relsize}
\usepackage{xspace}

% 使用 fontawesome（部分图标）
\usepackage{fontawesome} 

% A4纸，上下左右边距
\usepackage[
    a4paper,
    left=1.2cm,
    right=1.2cm,
    top=1.2cm,
    bottom=0.8cm,
    nohead
]{geometry}

\renewcommand{\baselinestretch}{1.3} % 行间距设为1.22

\usepackage{titlesec}
\usepackage{enumitem}
\setlist{noitemsep} % 取消列表项间的额外间距
%\setlist{nosep} % 取消所有垂直间距
\setlist[itemize]{topsep=0.25em, leftmargin=*}
\setlist[enumerate]{topsep=0.25em, leftmargin=*}

% --- 用于控制【不同项目之间】的垂直距离 ---
\newlength{\interProjectSpacing}
\setlength{\interProjectSpacing}{0.5em} % <--- 在此调整项目之间的距离
\newcommand{\projectsep}{\vspace{\interProjectSpacing}}

% --- 用于控制【项目标题】与下方【项目描述】的距离 ---
\newlength{\intraProjectTitleSep}
\setlength{\intraProjectTitleSep}{0.4em} % <--- 在此调整标题和描述的距离
\newcommand{\titlebreak}{\\[\intraProjectTitleSep]}

% --- 用于控制【项目描述】与下方【要点列表】的距离 ---
\newlength{\intraProjectListTopSep}
\setlength{\intraProjectListTopSep}{0.4em} % <--- 在此调整描述和列表的距离

% =======================================================================


\titleformat{\section}         % 定制 \section 命令 
{\large\bfseries\raggedright} % 将 section 标题设置为大号、粗体且左对齐
{}{0em}                      % 可用于为所有 section 添加前缀（如“章节...”）
{}                           % 可用于在标题前插入代码
[{\color{CVBlue}\titlerule}]  % 在标题后插入一条横线
\titlespacing*{\section}{0cm}{*1.6}{*1.2}



\begin{document}
\pagenumbering{gobble}

%%%% 利用tikz来定位照片
\begin{tikzpicture}[remember picture, overlay] 
    \node[anchor = north east] at ($(current page.north east)+(-2cm,-1.2cm)$) {\includegraphics[height=3cm]{photo.jpg}};
  \end{tikzpicture}%
  %%%% 利用tikz来定位学校Logo，这里只在第一页显示
  \begin{tikzpicture}[remember picture, overlay] 
    \node[anchor = north west] at ($(current page.north west)+(0.5cm,+1.0cm)$) {\includegraphics[height=6cm]{zju.png}};
  \end{tikzpicture}%
\centerline{\LARGE\bfseries{夏雨}} 

\begin{center}
\normalsize
\begin{tabular}{l @{\hspace{1cm}} l}  % 这里的 1cm 是两列之间的固定绝对间距
    \faPhone\ 158-6729-7122 & \faEnvelopeO\ \href{mailto:xy1@zju.edu.cn}{xy1@zju.edu.cn} \\
    \faHome\ 浙江·宁波 & \faCalendar\ 2001.10 \\
\end{tabular}
\end{center}


\section{\makebox[\widthof{\faGraduationCap}][c]{\color{CVBlue}\faGraduationCap}\ 教育背景}    

%%% ------------ 第一段经历：浙江大学 ------------ %%%
\noindent
\begin{tabularx}{\textwidth}{@{} p{6cm} @{\hspace{1.5em}} X @{\hspace{1em}} r @{}}
    \textbf{2024.09 - 至今} & \textbf{浙江大学 | 固态锂金属电池} & \hfill \textbf{硕士在读}
\end{tabularx}
\vspace{-2em} % 1. 这里改为负间距，强行把下方的列表往上拉

% 2. 增加 topsep=0pt 彻底清除 itemize 自带的顶部外边距
\begin{itemize}[nosep, topsep=0pt, leftmargin=1.5em]
    \item 主修课程：能源材料及电池技术、材料湿化学合成与制备、半导体物理。
    \item 荣誉奖项：浙江大学材料学院优秀团干
\end{itemize}

\vspace{0em} % 两段教育经历之间的间距

%%% ------------ 第二段经历：四川大学 ------------ %%%
\noindent
\begin{tabularx}{\textwidth}{@{} p{6cm} @{\hspace{1.5em}} X @{\hspace{1em}} r @{}}
    \textbf{2020.09 - 2024.06} & \textbf{四川大学 | 新能源材料与器件(A+)} & \hfill \textbf{学士}
\end{tabularx}
\vspace{-2em}
\begin{itemize}[nosep,topsep=0pt, leftmargin=1.5em]
    \item 主修课程：电工技术基础、电子技术基础、半导体物理、薄膜物理与技术、材料分析技术。
    \item 荣誉奖项：四川大学单项二等奖学金（连续三年）
\end{itemize}

\section{\makebox[\widthof{\faFileTextO}][c]{\color{CVBlue}\faFileTextO}\ 发表论文}

% ================== 论文/文献条目 ==================
\noindent 
\textbf{Mechanically Stable Honeycomb‐Like Gel Polymer Electrolyte Enabling Fast Li$^+$ Transport for High‐Performance Lithium Metal Batteries}\hfill\textbf{Advanced Functional Materials} \\ [-0.3em]
{\small \textbf{Author}：Hongwei Pan; Zhenyu Zhou; Yifei Li; Xueliang Zhang; \underline{Yu Xia}; Wei Guo; Sijie Xie; Pingping Tan; Xuan Zhang*; Jan Fransaer*}
\vspace{0.2em} % 如果有多篇文献，用这个控制文献条目之间的间距

% ================== 论文/文献条目 ==================
\noindent 
\textbf{Nitro functionalization and nanoscale confinement enable ether-based quasi-solid electrolytes with stable lithium metal and high-voltage compatibility}\hfill\textbf{Nanoscale} \\ [-0.3em]
{\small \textbf{Author}：Huiling Liu; Xingkai Jia; \underline{Yu Xia}; Zhu Liu*; Yinzhu Jiang*; Xuan Zhang*}% ==================================================

\vspace{0.2em} % 如果有多篇文献，用这个控制文献条目之间的间距

% ================== 论文/文献条目 ==================
\noindent 
\textbf{Tuning the Ion Transfer Behavior to Approaching Near‐Unity Li$^+$ Transference via Pore Engineering}\hfill\textbf{Small} \\ [-0.3em]
{\small \textbf{Author}：Xingkai Jia; Hongwei Pan; \underline{Yu Xia}; Pu Cheng; Shixiang Liu; Huiling Liu; Lingling Cao; Yinzhu Jiang*; Xuan Zhang*}% ==================================================
% 使用 \projectsep 命令来分隔两个项目
\projectsep

\section{\makebox[\widthof{\faFlask}][c]{\color{CVBlue}\faFlask}\ 科研经历}

% --- 第一个项目 ---
% 将标题行末尾的 \\ 替换为 \titlebreak 命令
\noindent
\begin{tabularx}{\textwidth}{@{} X @{\hspace{1em}} c @{\hspace{1em}} X @{}}
    \textbf{2024.11 - 2025.09} & \textbf{亚纳米MOF限域协同构筑快速锂离子动力学高性能固态电解质} & \hfill {负责人}
\end{tabularx}
\vspace{-2em} % 1. 这里改为负间距，强行把下方的列表往上拉

% 2. 增加 topsep=0pt 彻底清除 itemize 自带的顶部外边距
\begin{itemize}[nosep, topsep=0pt, leftmargin=1.5em]
    \item {\small \textbf{精通MOF材料的孔道环境调控，}能够通过几何限域与化学调控协同优化离子传输路径}
    \item {\small \textbf{掌握Zeta电位、XAFS、XPS等表征的测试与分析方法，}将实验与表征相结合，深入研究界面电化学环境}
    \item {\small \textbf{具备全电池组装能力，}独立制备电极材料与电解质，组装扣式与软包电池以评估不同构型下的材料性能}
\end{itemize}

\vspace{0em} 

% 使用 \projectsep 命令来分隔两个项目
\projectsep

% --- 第二个项目 ---
\noindent
\begin{tabularx}{\textwidth}{@{} l @{\hspace{1.5em}} X @{\hspace{1.5em}} r @{}}
    \textbf{2025.12 - 至今} & \textbf{基于额外金属离子引入的MOF路易斯酸位点调控及固态电解质性能提升} & \hfill {负责人}
\end{tabularx}
\vspace{-2em} % 1. 这里改为负间距，强行把下方的列表往上拉

% 2. 增加 topsep=0pt 彻底清除 itemize 自带的顶部外边距
\begin{itemize}[nosep, topsep=0pt, leftmargin=1.5em]
    \item {\small \textbf{掌握溶液浇筑法制膜工艺，}制备出兼具优异力学性能、显著阻燃性与良好柔韧性的复合固态电解质膜}
    \item {\small \textbf{掌握MOF后修饰及多金属调控技术，}能够构筑富路易斯酸位点功能填料并实现其结构与性能优化}
\end{itemize}

\section{\makebox[\widthof{\faUsers}][c]{\color{CVBlue}\faUsers}\ 项目经历}

% --- 第一个项目 ---
% 将标题行末尾的 \\ 替换为 \titlebreak 命令
\noindent
\begin{tabularx}{\textwidth}{@{} X @{\hspace{1em}} c @{\hspace{1em}} X @{}}
    \textbf{2024.11 - 2025.05} & \textbf{新型高比能长寿命正极材料储锂特性及失效机制研究(国家重点研发)} & \hfill {参与人}
\end{tabularx}
\vspace{-2em} % 1. 这里改为负间距，强行把下方的列表往上拉

% 2. 增加 topsep=0pt 彻底清除 itemize 自带的顶部外边距
\begin{itemize}[nosep, topsep=0pt, leftmargin=1.5em]
    \item {\small \textbf{了解卤化物固态电解质材料特性及其与电极界面的耦合机制，}能够通过键合调控促进可逆相变与氧化还原反应}
    \item {\small \textbf{熟悉高容量正极材料(如$\text{FeF}_3$、$\text{CuF}_2$)的失效机制与稳定策略，}具备界面化学修饰与CEI/SEI设计能力}
    \item {\small \textbf{精通电化学性能优化与衰减机制分析，}能系统评估循环稳定性、能量效率及其实际应用潜力}
\end{itemize}

\vspace{0em} 

% 使用 \projectsep 命令来分隔两个项目
\projectsep

% --- 第二个项目 ---
\noindent
\begin{tabularx}{\textwidth}{@{} l @{\hspace{1.5em}} X @{\hspace{1.5em}} r @{}}
    \textbf{2025.06 - 2026.04} & \textbf{正负极兼容多功能不燃电解液及耐高温高电压隔膜设计开发(国家重点研发)} & \hfill {参与人}
\end{tabularx}
\vspace{-2em} % 1. 这里改为负间距，强行把下方的列表往上拉

% 2. 增加 topsep=0pt 彻底清除 itemize 自带的顶部外边距
\begin{itemize}[nosep, topsep=0pt, leftmargin=1.5em]
    \item {\small \textbf{了解全电池安全性能评估标准与极端工况下衰减机制，}能够系统评估电池的不燃特性与高温高压热稳定性}
    \item {\small \textbf{熟悉高电压隔膜表面涂覆改性工艺与热失控阻隔方法，}具备多元复合隔膜构建与性能评测能力}
    \item {\small \textbf{精通多功能添加剂筛选与材料合成路径，}能通过功能基团设计匹配实际需求}
\end{itemize}


\section{\makebox[\widthof{\faCogs}][c]{\color{CVBlue}\faCogs}\ 专业技能}
\begin{itemize}[nosep]
    \item \textbf{软件技能：} 熟练掌握Origin数据绘图、Office办公、PS图像处理等
    \item \textbf{测试技能：} 熟练掌握电化学性能测试、XRD测试、SEM测试、Raman测试、FT-IR测试等
    \item \textbf{其他技能：} 英语CET-6
\end{itemize}
\end{document}